\documentclass[11pt,a4paper]{article}

\usepackage[utf8]{inputenc}
\usepackage[T1]{fontenc}
\usepackage{lmodern}
\usepackage{microtype}
\usepackage[english]{babel}
\usepackage[a4paper,margin=2.5cm,headheight=14pt]{geometry}
\usepackage{xcolor}
\usepackage{listings}
\usepackage{enumitem}
\usepackage{booktabs}
\usepackage{longtable}
\usepackage{fancyhdr}
\usepackage{hyperref}
\usepackage{titlesec}
\usepackage{tcolorbox}

\hypersetup{
    colorlinks=true,
    linkcolor=black,
    urlcolor=blue!60!black,
    citecolor=black,
}

\definecolor{codebg}{RGB}{248,248,242}
\definecolor{codecomment}{RGB}{106,135,89}
\definecolor{codekeyword}{RGB}{0,102,153}
\definecolor{codestring}{RGB}{170,55,55}
\definecolor{ackcolor}{RGB}{170,90,30}

\lstdefinestyle{recipix}{
    backgroundcolor=\color{codebg},
    basicstyle=\ttfamily\small,
    keywordstyle=\color{codekeyword}\bfseries,
    commentstyle=\color{codecomment}\itshape,
    stringstyle=\color{codestring},
    breaklines=true,
    showstringspaces=false,
    frame=single,
    framerule=0pt,
    framesep=4pt,
    aboveskip=0.4em,
    belowskip=0.4em,
    morekeywords={recipe,function,ingredient,step,let,if,else,
        repeat,times,foreach,in,at,for,evaluate,scale,substitute,
        serves,with,by,ratio,return,quantity\_of,true,false,
        int,float,bool,Mass,Volume,Count,Temperature,Duration,Pinch,
        combine,mix,pour,melt,whisk,blend,bake,flip,add,sprinkle,drizzle,knead},
    morecomment=[l]{//},
    morestring=[b]",
}

\lstdefinestyle{output}{
    backgroundcolor=\color{codebg},
    basicstyle=\ttfamily\small,
    breaklines=true,
    showstringspaces=false,
    frame=single,
    framerule=0pt,
    framesep=4pt,
    aboveskip=0.4em,
    belowskip=0.4em,
}

\lstdefinestyle{py}{
    backgroundcolor=\color{codebg},
    basicstyle=\ttfamily\footnotesize,
    keywordstyle=\color{codekeyword}\bfseries,
    commentstyle=\color{codecomment}\itshape,
    stringstyle=\color{codestring},
    breaklines=true,
    showstringspaces=false,
    frame=single,
    framerule=0pt,
    framesep=4pt,
    aboveskip=0.4em,
    belowskip=0.4em,
    language=Python,
}

\lstset{style=recipix}

\titleformat{\section}{\Large\bfseries}{\thesection.}{0.6em}{}
\titleformat{\subsection}{\large\bfseries}{\thesubsection}{0.6em}{}
\titleformat{\subsubsection}{\normalsize\bfseries}{}{0em}{}

\pagestyle{fancy}
\fancyhf{}
\fancyhead[L]{CSE 341 -- Recipix v4.1}
\fancyhead[R]{D4 (Part 2 continuation)}
\fancyfoot[C]{\thepage}

\newcommand{\entryheader}[4]{%
    \par\vspace{1em}\noindent\textbf{Entry \#{#1}} \textbullet{} \textbf{Date} #2 \textbullet{} \textbf{Phase} #3 \textbullet{} \textbf{AI tool} #4 \par\vspace{0.3em}\noindent\rule{\textwidth}{0.4pt}\par\vspace{0.3em}}

\newcommand{\sessionheader}[1]{%
    \par\vspace{1.5em}\noindent\textbf{Source session: }\texttt{#1}\par\vspace{0.5em}}

\newcommand{\field}[1]{\noindent\textbf{#1}}

\title{\textbf{D4 -- AI Usage Journal (Part 2 continuation)} \\
       \large Recipix v4.1 -- CSE 341 Semester Project}
\author{Berk Hakan \"Oge (210104004132)}
\date{23 May 2026}

\begin{document}
\maketitle
\thispagestyle{fancy}


\section*{Continuation note}

This document continues the AI Usage Journal I submitted with Part~1
(\texttt{210104004132\_D4\_P1.pdf}, 15 entries across 6 sessions, dated
May 6--8 2026). It documents the AI sessions I used during Part~2 work,
dated May~20--23 2026. Entry numbering continues from the Part~1
journal: this document contains entries \textbf{16 through 24}.

Part~2 added two AI workflows on top of the Part-1 pattern (web
brainstorming + PM-in-Code + autonomous developer/tester agents).
First, a single Claude.ai web planning chat that turned into the
Part~2 roadmap and the \texttt{part2\_plan.md} contract with İsmail.
Second, three discrete Claude.ai web sessions running the mandatory
experiments E1 (catch-up -- see entry~17 note), E2, and E3 each in a
fresh model with no prior context. Third, three Claude Code sessions
covering the Part~2 build: a developer agent that produced the
interpreter + 36-test suite, a PM agent that reviewed those commits
and closed three gaps, and a professor agent that ran a Pass-1
rubric-graded review of the in-scope artifacts. The professor session
was extended on May~22--23 with a second turn (final rubric review of
the finalized state -- same session UUID, two entries cite it), and
a fresh tester session was run on May~23 for the final functional
verification before submission.

In total Part~2 produced \textbf{8 sessions} (sessions 7--14 in the
table below) and \textbf{9 entries} (entries 16--24), since the
professor session is cited by two entries.

\subsection*{Sessions covered in this continuation}

\begin{center}
\small
\begin{tabular}{@{}clll@{}}
\toprule
\textbf{\#} & \textbf{Session} & \textbf{Tool / Model} & \textbf{Date} \\
\midrule
7  & Web Planning B -- Part 2 Roadmap and Pair-Split            & Claude.ai web (Sonnet)                        & 2026-05-20 \\
8  & Web Experiment E1 -- EBNF generation (late catch-up)       & Claude.ai web (Opus 4.7, thinking on)          & 2026-05-21 \\
9  & Web Experiment E2 -- Type equivalence                      & Claude.ai web (Opus 4.7, thinking on)          & 2026-05-21 \\
10 & Web Experiment E3 -- AI implements \texttt{\_check\_ActionStmt} & Claude.ai web (Opus 4.7, thinking on)      & 2026-05-21 \\
11 & Developer Agent -- Interpreter + Rendering build           & Claude Code (Sonnet 4.6, autonomous)          & 2026-05-21 \\
12 & PM Agent -- Interpreter Commits Review                     & Claude Code (Opus 4.7)                        & 2026-05-21 \\
13 & Professor Agent -- Pass-1 + Final Rubric Review            & Claude Code (Opus 4.7)                        & 2026-05-21 / 2026-05-23 \\
14 & Tester Agent -- Final Pre-Submission Verification          & Claude Code (Opus 4.7)                        & 2026-05-23 \\
\bottomrule
\end{tabular}
\end{center}

The three Claude Code transcripts are submitted alongside this PDF as
JSONL files:

\begin{itemize}[leftmargin=*,itemsep=0pt]
\item \texttt{session11\_developer\_interpreter\_2026-05-21.jsonl}
\item \texttt{session12\_pm\_interpreter\_review\_2026-05-21.jsonl}
\item \texttt{session13\_professor\_2026-05-21\_to\_2026-05-23.jsonl}
\item \texttt{session14\_tester\_final\_verification\_2026-05-23.jsonl}
\end{itemize}

The four web sessions have shareable URLs cited inline within each
entry.


\section{Session 7 -- Web Planning B: Part 2 Roadmap (Claude.ai Web)}

Source session: full transcript saved to repo as \texttt{web\_ses.txt}
at submission time. Single Claude.ai web chat used to plan the Part~2
split with my partner, generate the contract proposal doc to send to
İsmail, and surface the design decisions that needed locking before
either of us wrote new code. Four substantive turns, treated as one
entry because they form one continuous planning arc.

\entryheader{16 of 24}{2026-05-20}{Project planning / pair split}{Claude.ai web (Sonnet)}

\field{Goal.} Take Part 1's submitted state into a Part 2 plan İsmail
and I could both sign off on before either of us started writing code.
Specifically: pick a pair split, lock the integration contracts that
would let us work in parallel, and produce a markdown doc I could
send to İsmail rather than walking him through chat scrollback.

\field{Prompt.} Started with \emph{"We are building our own programming
language. Here the github repo for it. We have completed Part1 of the
project. You can read the project handout.pdf file for details about
the project. Read trough the files and inspect the work done overall
for the part1. After you are done with learning the project we will
start planning on what to do next for part2."} Then three follow-up
turns: ask for split options, pick Option~A, ask for the markdown doc
to share with İsmail.

\field{Response (key part).} The model pulled the spec, parser source,
AST, error module, lexer contract, fixtures, my D5, İsmail's D5, the
D6, and the handout. Three things it produced were load-bearing: (1)~a
clean read of where Part~1 actually shipped, including the three
deferred items I'd flagged in my D5 (\texttt{AmbiguousCall} resolution,
\texttt{evaluate}-only constraint, single-return rule) and which Part~2
phase they each landed in. (2)~Three pair-split options, each with a
Sebesta-chapter argument; recommended Option~A (Berk = interpreter,
İsmail = type checker) was justified by İsmail's stated D5 intention
and by my ownership of the AST shapes. (3)~A full Phase-0 contract list
(TypedAST contract, runtime value dataclasses, action-verb signature
table, quantity-normalization rule, environment-chain shape, file
layout, sample-output agreement) plus a draft \texttt{part2\_plan.md}
written in my voice.

\field{Accepted.} The Option-A pair split; the in-place TypedAST
annotation contract; the runtime value dataclasses as drafted
(\texttt{RtQuantity}, \texttt{RtPinch}, \texttt{RtIngredient},
\texttt{RtRecipe}, \texttt{RtList}); the quantity-normalization rule
(type checker tracks dimension, interpreter normalizes to base units
at \texttt{QuantityLit} evaluation); the shared \texttt{environment.py}
design. The first draft of \texttt{part2\_plan.md} became the working
contract; subsequent revisions (Rev~2, Rev~3) were edits in a separate
Claude Code session, but the bones came from here.

\field{Rejected / modified.} The 12-row per-verb action-verb signature
table was the model's straw proposal; I collapsed it to three signature
classes (Heterogeneous, Homogeneous, Nullary) in
\texttt{part2\_plan.md}~\S2.3 because the per-verb form was unjustified
branch explosion in the checker. The model accepted the collapse
without argument when challenged later -- see Session~10 / Entry~19.

\field{Errors caught.} One soft error: the model's first read called
the action-verb signature table ``small but exam-bait'' and proposed
twelve rows. The right answer was three. I caught it by reading the
\S2.3 list against the spec's decision~\#29 (``action verbs are an
exception to the homogeneous-list rule'') -- twelve per-verb
signatures don't honor a single-decision carve-out; three signature
classes do. The model's later proposal had the same blind spot. This
is the failure mode I now look for first when prompting any AI on
design decisions: it expands tables in the direction of ``more rules
cover more cases'' rather than collapsing them in the direction of
``fewer rules cover the same cases at lower exam and maintenance
cost.''

\field{Reflection.} This session's contribution was that it forced me
to articulate the Part~2 split out loud before talking to İsmail. The
chapter mapping (Ch.~6 = İsmail, Ch.~7 = Berk) is the line I rehearsed
for the meeting. The runtime contract doc is what made the parallel
build possible the next day. The model was best at
\emph{briefing-from-spec} (it correctly identified the three deferred
items from my Part~1 D5 without me prompting), worst at design
rationale that has to be in my own voice -- every D1~[P2] paragraph
it offered, I rewrote. Lesson: I use the AI to surface decisions I
should make, never to make them.


\section{Session 8 -- Web Experiment E1: EBNF Generation}

\begin{tcolorbox}[colback=ackcolor!8,colframe=ackcolor!50,boxrule=0.5pt]
\small\textbf{Late-execution acknowledgment.} Experiment E1 was
required for the Part~1 D4 journal but was omitted from my
submission by oversight. I noticed the gap on 2026-05-21 during a
Part~2 D4 prep pass and am running E1 here, honestly dated, as a
catch-up entry filed with my Part~2 D4 submission. This entry does
\textbf{not} retroactively satisfy the Part~1 requirement; it
documents the experiment for completeness and is noted in my D5
retrospective. The model used (Claude Opus~4.7) postdates Part~1's
working window, which is another reason to be explicit about
dating rather than smuggle it into earlier dates.
\end{tcolorbox}

\sessionheader{claude.ai/share/6d67cab6-006e-414f-b05c-ada5fa5b0d22}

Two-turn session: cold EBNF generation, then comparison against the
locked spec.

\entryheader{17 of 24}{2026-05-21}{Design / EBNF retrospective}{Claude Opus 4.7 (extended thinking: on)}

\field{Goal.} Generate a first-pass EBNF grammar for Recipix from a
short prose description and compare it against my locked spec
(v4.1, \S13), so I can see which language-specific decisions an AI
derives on its own and which it silently smooths over into PL-textbook
defaults.

\field{Prompt.} Verbatim, one paragraph -- no productions or repository
content leaked into turn~1; only the one-paragraph description of the
language: ``Recipix is a small DSL for writing parameterized recipes
that can be scaled and have ingredients substituted. The primitive
types are \texttt{int}, \texttt{float}, \texttt{bool}, five quantity
types (\texttt{Mass}, \texttt{Volume}, \texttt{Count},
\texttt{Temperature}, \texttt{Duration}), and a separate \texttt{Pinch}
type$\dots$'' -- continues through structured types, user-visible
constructs, the three call-site operators, quantity literals, the
unary \texttt{quantity\_of} operator, and the no-chained-comparison
constraint. Ends with ``Please give me an EBNF grammar.''

\field{Response (key part, turn~1).} AI's first-pass EBNF, generated
cold. Highlights:

\begin{lstlisting}
Expression     ::= LogicalOr
LogicalOr      ::= LogicalAnd ( "or"  LogicalAnd )*
LogicalAnd     ::= Comparison ( "and" Comparison )*
Comparison     ::= Additive ( CompOp Additive )?       (* non-associative *)
CompOp         ::= "<" | "<=" | ">" | ">=" | "==" | "!="
Additive       ::= Multiplicative ( ( "+" | "-" ) Multiplicative )*
Multiplicative ::= Unary ( ( "*" | "/" | "%" ) Unary )*
Unary          ::= ( "-" | "not" | "quantity_of" | "evaluate" ) Unary
                 | Postfix
SubstituteExpr ::= "substitute" "(" Expression "," Expression ","
                                   "with"  ":" Expression ","
                                   "ratio" ":" Expression ")"
QuantityLit    ::= ( INT_LIT | FLOAT_LIT ) Unit
\end{lstlisting}

\noindent Commentary from the AI: quantity literals as \emph{two tokens}
``joined in the parser''; chained comparison ``forbidden by making
Comparison non-associative''; \texttt{if}/\texttt{else} ``no
dangling-else problem because bodies are braced''; \texttt{quantity\_of}
``sits at the unary layer''. The AI volunteered uncertainty about
\texttt{Count} (``I left it without a dedicated unit on the assumption
it's promoted from INT\_LIT by context'') and flagged it as the most
likely divergence.

\field{Response (key part, turn~2).} After giving the AI the GitHub
link to \texttt{docs/recipix\_v4\_1\_spec.md}, its self-assessment:
\emph{``I got the high-stakes structural calls right -- two-token
quantities, non-associative comparison, mandatory if/else braces,
quantity\_of at the unary level -- but missed the substitute
bare-identifier restriction, and several smaller surface details.''}

\begin{center}
\small
\begin{tabular}{@{}p{0.32\textwidth}p{0.30\textwidth}p{0.30\textwidth}@{}}
\toprule
\textbf{Decision} & \textbf{AI cold} & \textbf{Locked spec} \\
\midrule
\#7 quantity literals as two tokens & two tokens, parser-joined & match (AI missed \texttt{count}, mg, used wrong unit set) \\
\#17 comparison non-associative & non-associative & match, same reasoning \\
\#20 mandatory braces on if/else & mandatory & match \\
\#34 \texttt{quantity\_of} at unary & unary level, parens-free operand & ``right floor, wrong storefront'' -- match on placement \\
\#35 substitute slots as IDENT & all four slots \texttt{Expression} & \textbf{miss} -- the one big design call the AI invented from prose was wrong \\
\bottomrule
\end{tabular}
\end{center}

The AI's framing of the substitute miss: \emph{``Your AST node carries
original\_name: str and replacement\_name: str, not Expr -- there is no
AST shape that pretends the slot could be computed. That's the cleanest
design lesson in the whole comparison -- grammar-as-documentation-of-semantics,
not just grammar-as-parser-spec.''}

\field{Accepted.} The four high-stakes calls match the locked grammar's
intent, and the AI's reasoning on each lines up with the rationale I
locked into the spec. The precedence ladder shape (seven levels in the
right top-down order) and the program skeleton both match. Even the
``right floor, wrong storefront'' miss on \texttt{quantity\_of} proves
the type-system reasoning was sound; only the surface syntax differed.

\field{Rejected / modified.} The substitute-slot miss plus a cluster of
over-permits where the spec locks down: optional \texttt{let} type
annotation, \texttt{List<T>} as a user-writable type, top-level
\texttt{IngredientDecl}, step body absent, no closed action-verb set,
function bodies with no single-return restriction, postfix indexing,
keyword logical operators (\texttt{and}/\texttt{or}/\texttt{not}
instead of \texttt{\&\&}/\texttt{||}/\texttt{!}), wrong unit set.

\field{Errors caught.} (1)~\textbf{Substitute slots as \texttt{Expression}
instead of \texttt{IDENT}} -- under the AI grammar,
\texttt{substitute(r, flour + salt, with: milk, ratio: 1.0)} parses;
under the spec it's a parse error. (2)~\textbf{Cluster of over-permits}
where the spec locks down -- predictable shape of ``derive a grammar
from prose without the locked decisions in view.'' (3)~\textbf{Unit set
drift} -- AI invented units (\texttt{°F}, \texttt{s}, \texttt{h}) and
missed locked ones (\texttt{count}, \texttt{mg}).

\field{Reflection.} The thing that surprised me here is how much Opus
got right cold with thinking on. Four of five big design calls came
back match. What didn't change is the \emph{kind} of thing it still
misses. Decision~\#28 (``ingredient identity is the binding'') drives
decision~\#35 (substitute slots are bare identifiers), and there's no
way to derive \#35 from prose unless you've already had the \#28
conversation. The over-permits -- \texttt{List<T>} annotations as
user-writable, top-level ingredients, optional \texttt{let} type,
indexing, multi-return function bodies -- are the same shape: the AI
admits whatever generic PL languages admit; the spec restricts what
Recipix specifically restricts. \emph{Lesson:} I can't trust a
``generate a grammar from a paragraph'' pass to surface decisions
whose rationale lives outside that paragraph. I have to drive the AI
with the rationale already in hand. That's also why I'm filing this
entry on its real date and not backdated -- the value of the
experiment is in the comparison, and the comparison only holds if
the dates and tool name are honest.


\section{Session 9 -- Web Experiment E2: Type Equivalence}

\sessionheader{claude.ai/share/dd9a44cf-584d-4194-9312-2b512752f90d}

Two-turn session: cold equivalence recommendation, then comparison
against the locked spec.

\entryheader{18 of 24}{2026-05-21}{Type checker / Design}{Claude Opus 4.7 (extended thinking: on)}

\field{Goal.} Test whether a fresh AI session, given only a
one-paragraph description of Recipix and its structured types, can
derive the type-equivalence rule I actually settled on in v4.1 -- or
whether it gives a textbook answer that misses the domain reason for
the split.

\field{Prompt.} Verbatim, Recipix domain framing plus the four
structured types (\texttt{Ingredient}, \texttt{Recipe}, \texttt{Step},
\texttt{List<T>}), asking for a per-type name-vs-structural
recommendation plus a small example program that depends on the rule.

\field{Response (key part, turn~1, AI's first-pass recommendation):}

\begin{center}
\small
\begin{tabular}{@{}lp{0.20\textwidth}p{0.55\textwidth}@{}}
\toprule
\textbf{Type} & \textbf{AI cold} & \textbf{AI's justification (paraphrased)} \\
\midrule
\textbf{Ingredient} & name & ``Domain noun... generic fields would collapse with hypothetical Cost/Allergen records.'' \\
\textbf{Recipe} & name & ``Top-level entity; accidental equivalence at the top of the type lattice is a worse failure mode than at the leaves.'' \\
\textbf{Step} & name & ``Optional fields make decisive -- under structural rules you'd get four distinct shapes per optional-field combination.'' \\
\textbf{List\textless T\textgreater} & structural in parameter, nominal in constructor & ``Standard treatment for parametric types -- ML/Haskell/Java generics converge here.'' \\
\bottomrule
\end{tabular}
\end{center}

The AI offered a brownies-recipe example using pseudo-syntax
(\texttt{Step("desc", actions: [...])} constructor calls inside a
\texttt{steps: [...]} list literal) claiming the three Steps in the
brownies recipe (one with neither modifier, one with both, one with
only \texttt{duration}) would have ``three distinct structural
shapes'' and fail to unify in a \texttt{List<Step>} under structural
equivalence.

\field{Response (key part, turn~2, comparison against the locked spec):}
Self-assessment:

\begin{enumerate}[leftmargin=*,itemsep=2pt]
  \item \textbf{Recipe -- match.} Same call, same reasoning.
  \item \textbf{List\textless T\textgreater\ -- match.} Same conceptual
        move, phrased differently.
  \item \textbf{Ingredient -- divergence, with self-correction.}
        \emph{``My reasoning was forward-looking: avoid silent
        collisions with hypothetical Cost \{name, quantity\} or
        Allergen \{name, quantity\} types in some v2. That argument
        falls apart against your actual v1, because there is no
        user-defined record syntax -- Ingredient is the only \{name,
        quantity\} record the type system will ever see, and \S1 makes
        the type-name set closed. So my choice was buying safety
        against a threat that can't exist in the language as specified.
        Pick the more permissive rule when over-restriction has no
        upside.''}
  \item \textbf{Step -- divergence on reasoning.} The AI realized its
        argument was based on the strict-Sebesta reading of structural
        equivalence (each \texttt{?}-field present-or-absent creates a
        distinct shape) -- \emph{not} how the spec applies the rule.
        The spec treats optional fields as part of the \emph{single}
        declared shape.
\end{enumerate}

The AI's one-line summary: \emph{``I got Ingredient genuinely wrong
(defensible reasoning, wrong call for a sealed v1), I got Step
nominally wrong but on a question your spec doesn't actually need to
answer, and my example was sold as a discriminating test case when it
isn't one under your rules.''}

\field{Accepted.} Recipe and List\textless T\textgreater\ calls match
the spec's intent. AI's framing in turn~2 -- ``pick the more permissive
rule when over-restriction has no upside'' -- is a one-liner defense of
decision~\#10's structural-for-Ingredient call that I can quote almost
verbatim in D1~\S4.8.

\field{Rejected / modified.} Three substantive divergences:
(1)~Ingredient picked name (spec picks structural).
(2)~Step reasoning relied on strict-Sebesta reading; spec uses
optional-fields-inclusive reading. (3)~Pseudo-syntax in the example
(constructor-call form, list-literal of steps) is not Recipix surface
syntax.

\field{Errors caught.} (1)~The AI's domain reasoning optimized for
extensibility scenarios that the closed v1 type-name set forecloses.
(2)~The AI assumed a stricter notion of structural equivalence than
the spec uses. (3)~The example program does not run under the rule it
claimed to demonstrate -- the literal test the handout's E2
specification names.

\field{Reflection.} The interesting result here isn't that the AI got
Ingredient wrong -- it's that the reasoning behind the wrong answer
was sophisticated and forward-looking, and would have been correct for
a different language. It hedged against \texttt{Cost \{name, quantity\}}
collisions that v1 makes impossible because the type-name set is
closed. That's the failure mode I think I'll see most often when
prompting AI on Recipix design questions: the AI optimizes for the
kind of evolution a generic record-based DSL would undergo, and v1
has already foreclosed those evolutions. \emph{Lesson:} the AI
accelerates generating design options. It does not decide which
option a sealed v1 has already committed to. That's mine.


\section{Session 10 -- Web Experiment E3: AI Implements \texttt{\_check\_ActionStmt}}

\sessionheader{claude.ai/share/8190d922-7330-4181-84cd-40e1748c1e99}

Two-turn session: cold implementation, then comparison against plan
\S2.3 and the typechecker stub.

\entryheader{19 of 24}{2026-05-21}{Type checker}{Claude Opus 4.7 (extended thinking: on)}

\field{Goal.} Have the AI implement \texttt{\_check\_ActionStmt} on the
\texttt{TypeChecker} class -- the \S10 error-\#13 enforcer for the
twelve closed verbs in spec~\S8 -- without first showing it the
three-class table from plan~\S2.3, then compare its inferred
signatures to the locked design.

\field{Prompt.} Verbatim: the twelve verbs (\texttt{combine, mix, add,
sprinkle, pour, drizzle, whisk, blend, knead, melt, bake, flip}), the
type-string contract for \texttt{env.lookup}, the four test inputs
whose right answers I knew in advance, and the requirement to use at
least one helper data structure (a verb-signature dict/table).

The AI flagged an unsolicited prompt ambiguity before answering:
\emph{``you said D $\in$ \{Mass, Volume, Count, Temperature, Duration\}
-- Pinch is not in D -- but then the example return value mentions
Quantity\textless Pinch\textgreater, and test \#1 uses
Ingredient\textless Pinch\textgreater''}. Real prompt bug; the spec
treats Pinch as a primitive, not a dimension, but the implicit-projection
rule (\#31) and the test inputs effectively require the checker to
handle \texttt{Ingredient<Pinch>} as well-formed.

\field{Response (key part, turn~1, AI's first-pass implementation):}

\begin{lstlisting}[style=py]
VERB_SIGS = {
    'combine': ('variadic', [], _ingredient_of(), 2),
    'mix':     ('variadic', [], _ingredient_of(), 2),
    'blend':   ('variadic', [], _ingredient_of(), 2),
    'whisk':   ('fixed',    [_ingredient_of()]),
    'flip':    ('fixed',    [_ingredient_of()]),
    'knead':   ('fixed',    [_ingredient_of('Mass')]),
    'melt':    ('fixed',    [_ingredient_of('Mass')]),
    'pour':    ('fixed',    [_ingredient_of('Volume')]),
    'drizzle': ('fixed',    [_ingredient_of('Volume')]),
    'sprinkle':('fixed',    [_ingredient_of('Mass', 'Count')]),
    'add':     ('variadic', [_ingredient_of()], _quantity_of(), 1),
    'bake':    ('fixed',    [_ingredient_of(),
                             _quantity_of('Temperature'),
                             _quantity_of('Duration')]),
}
\end{lstlisting}

Dispatcher distinguishes two signature kinds (\texttt{'fixed'} and
\texttt{'variadic'}) with per-arg shape-check delegated to predicate
factories. The AI explicitly named its biggest guesses up front:
\emph{``combine/mix/blend accept any combination of dimensions...
bake is fixed-arity 3. This is the biggest guess in the table.''}

\field{Response (key part, turn~1, predicted verdicts on the four tests):}

\begin{center}
\small
\begin{tabular}{@{}clll@{}}
\toprule
\textbf{\#} & \textbf{Test input} & \textbf{AI verdict} & \textbf{AI reason} \\
\midrule
1 & \texttt{combine(flour, salt)}        & accept & variadic $\geq 2$, both Ingredients \\
2 & \texttt{pour(flour:Mass)}            & reject & pour locked to Volume \\
3 & \texttt{bake(x)}                     & reject & bake arity 3, called with 1 \\
4 & \texttt{blend(milk:Vol, sugar:Mass)} & accept & blend in any-mixed-dim bucket \\
\bottomrule
\end{tabular}
\end{center}

\field{Response (key part, turn~2, comparison against plan~\S2.3):}
\emph{``The verdict is uncomfortable but useful for your E3 entry --
I diverged in three substantive ways, and two of four test predictions
flip under \S2.3.''} The three substantive divergences (AI's words):

\begin{enumerate}[leftmargin=*,itemsep=2pt]
  \item \textbf{Per-verb table vs three-class collapse.}
        \emph{``I built a per-verb table. \S2.3 collapses twelve verbs
        into three classes precisely to avoid the per-verb branch
        explosion. I optimized for cooking semantics; you optimized for
        type-system parsimony.''}
  \item \textbf{Heterogeneous-group misses.} Partial credit on
        \texttt{combine}/\texttt{mix} (required $\geq 2$ args instead of
        $\geq 1$); \texttt{add} was bespoke (spec puts it in
        heterogeneous); \texttt{sprinkle} was restricted to Mass/Count
        (spec admits any dimension including Pinch).
  \item \textbf{Homogeneous-group misses.} Most divergent. Locked
        per-verb dimensions; missed homogeneity rule entirely; missed
        \texttt{flip} as nullary. Also missed decision~\#31.
\end{enumerate}

Four-test re-run under \S2.3:

\begin{center}
\small
\begin{tabular}{@{}clllc@{}}
\toprule
\textbf{\#} & \textbf{Test} & \textbf{AI cold} & \textbf{\S2.3} & \textbf{Status} \\
\midrule
1 & \texttt{combine(flour, salt)}     & accept & accept              & match (right verdict, wrong rule) \\
2 & \texttt{pour(flour:Mass)}         & reject & \textbf{accept}     & \textbf{flipped} \\
3 & \texttt{bake(x)}                  & reject (arity 3$\neq$1) & reject (arity 0$\neq$1) & match-by-coincidence \\
4 & \texttt{blend(milk, sugar)}       & accept & \textbf{reject}     & \textbf{flipped} \\
\bottomrule
\end{tabular}
\end{center}

AI's score framing: \emph{``If you're scoring strictly on `did the AI
produce the same accept/reject column as the spec,' 2/4. If you're
scoring on `did the AI produce the right verdict for the right rule,'
1/4.''} The AI then volunteered a \S2.3-aligned correction unsolicited
-- a clean three-class dispatcher with sets + \texttt{\_project\_to\_quantity}
helper. That correction became a usable starting point for the
implementation.

\field{Accepted.} Dispatcher skeleton is reusable. The
\texttt{\_parse\_type} helper survives unchanged. AI's
\texttt{error\_code=13} choice is consistent with the spec. The AI's
corrected \texttt{\_check\_ActionStmt} body in turn~2 became part of
the eventual implementation. AI's unsolicited prompt-ambiguity flag
(Pinch in D vs not in D) surfaced something the v4.2 spec needs to
clarify.

\field{Rejected / modified.} Everything in the cold first-pass except
the dispatcher skeleton: the twelve-row \texttt{VERB\_SIGS} table
(replaced with the three-class form), per-verb dimension locks (no
per-verb dim in spec, only homogeneity), \texttt{combine}/\texttt{mix}/\texttt{blend}
min-arity $\geq 2$ (spec is $\geq 1$), bespoke \texttt{add} signature
(spec heterogeneous), \texttt{bake} as fixed-arity 3 (spec nullary;
temperature/duration moved to step modifiers), \texttt{flip} unary
(spec nullary), \texttt{whisk} unary (spec homogeneous), missing
implicit Ingredient~$\to$~Quantity projection.

\field{Errors caught.} (1)~Test 2 (\texttt{pour(flour)}) verdict flipped
-- AI rejected on hard-coded Volume; spec accepts because pour is
homogeneous and single-arg is trivially homogeneous.
(2)~Test 4 (\texttt{blend}) verdict flipped -- AI grouped with combine/mix
as heterogeneous; spec puts in homogeneous with dimension mismatch.
(3)~Test 3 matched verdict for the wrong reason.
(4)~Decision~\#31 (Ingredient~$\to$~Quantity projection) absent.
(5)~Per-verb cooking-semantics over-specify against the spec's
type-system parsimony.

\field{Reflection.} The biggest takeaway from this one is that the
AI's design instinct -- ``more rules cover more cases'' -- is exactly
opposite to the spec's instinct of ``fewer rules cover the same cases
at lower exam and maintenance cost.'' Twelve per-verb signatures look
like authority on first read; the AI's own retrospective lands on this
in one line: \emph{``I optimized for cooking semantics; you optimized
for type-system parsimony.''} That framing is the one I want in
\S4.8 and in my head for the exam. The other half of the experiment
is what happened after the spec came into view: the AI wrote its own
\S2.3-aligned correction, unsolicited, and got it essentially right.
\textbf{AI implements decisions well. It does not invent them well.}
That's the working pattern I'll keep across the rest of Part~2 -- I
drive the decisions, the AI drives the implementation, and the
journal is the record of what happens when those roles slip.


\section{Session 11 -- Developer Agent: Interpreter + Rendering Build}

Source: \texttt{session11\_developer\_interpreter\_2026-05-21.jsonl}
(uploaded alongside this PDF). Autonomous developer-agent run
dispatched from a Claude Code PM session, executed in a single long
autonomous turn.

\entryheader{20 of 24}{2026-05-21}{Interpreter implementation (autonomous agent)}{Claude Sonnet 4.6 (Claude Code agent)}

\field{Goal.} Build the runtime half of Recipix Part~2:
\texttt{interpreter.py} and \texttt{rendering.py}, plus a 36-test suite
covering the \S4 task list in \texttt{part2\_plan.md}. The type checker
remains stubbed (İsmail's half); the agent must not dispatch
\texttt{AmbiguousCall} and must fail loudly if it sees one. Sample~1
and Sample~2 must produce rendered cookbook-style output; Sample~3
must remain a type-error path.

\field{Prompt.} \textasciitilde 3.5k-character developer-agent prompt
covering: scope (interpreter + rendering only; do not modify Part~1
files; do not touch \texttt{typechecker.py}), the runtime contract
sections of \texttt{part2\_plan.md} to read first, the 14-step
recommended order, the runtime-error checklist (5 errors from spec
\S10), and the acceptance criteria.

\field{Accepted.} The agent produced commits \texttt{51c1b03}
(interpreter + rendering implementation, 619 lines) and \texttt{fac696a}
(36-test suite). Tools used: $\sim$40 Edit + Write, $\sim$30 Read,
$\sim$20 Bash for verification. The eager-step-evaluation strategy was
the one locked at the time in plan~\S2.2; the agent followed it.
\texttt{AmbiguousCall} dispatch raises \texttt{RuntimeRecipixError}
unconditionally -- the loud-fail behavior asked for; integration tests
use a separate helper to simulate the typechecker's resolution.

\field{Rejected / modified.} Three findings from a follow-up PM review
(entry~21) led to commit \texttt{d5c539c}: the original AmbiguousCall
handler had a fallback dispatching via symbol-table \texttt{isinstance}
checks (would let a missing-typechecker bug pass silently for any name
in scope), and the test suite had gaps on runtime error~\#4
(\texttt{substitute} of unknown ingredient) and decision~\#31. All
three were closed in the follow-up commit; the agent self-corrected on
all three when given the gap report.

\field{Errors caught.} The PM review (entry~21) caught all three gaps
before any of them shipped to a passing-on-the-surface-but-broken-in-integration
state. The agent itself self-caught only one error during the build:
an early version of the unit-conversion table had a wrong conversion
factor; it noticed during the test pass for sample~1 and corrected.
Test count: 36, all passing on first full run.

\field{Reflection.} The pattern I keep seeing with autonomous developer
agents in Claude Code is that when the contract is locked tightly
enough, the agent's free movement is on the implementation surface
only -- it doesn't drift into design decisions, it doesn't invent
abstractions I didn't ask for, and it ticks the task list. \emph{AI
is good at breadth-first implementation when the spec is locked} --
exactly the same finding as my Part~1 entries 14 and 15, and the same
caveat applies: AI is unreliable on second-order semantics. The
eager-vs-lazy step-eval question is one such second-order question
the agent didn't surface; I had to lock it in the contract before the
build and the agent followed.


\section{Session 12 -- PM Agent: Interpreter Commits Review}

Source: \texttt{session12\_pm\_interpreter\_review\_2026-05-21.jsonl}.
PM-agent session that reviewed the two developer commits against the
runtime contract in \texttt{part2\_plan.md}.

\entryheader{21 of 24}{2026-05-21}{Post-developer review (PM agent)}{Claude Opus 4.7}

\field{Goal.} Have the PM read the developer's interpreter + test
commits, run the test suite, and produce a realistic review (not a
rubber stamp) before shipping the work to İsmail's typechecker handoff.

\field{Prompt.} PM-agent prompt covering: scope (commits \texttt{51c1b03}
and \texttt{fac696a} only; do not modify code; do not touch
\texttt{typechecker.py}; produce a written review with file:line
citations), the contract sections to verify against, the six rubric
dimensions, and a structured deliverable.

\field{Accepted.} PM produced a structured review with per-contract
verdicts. Three gaps surfaced:

\begin{itemize}[leftmargin=*,itemsep=2pt]
  \item \textbf{Gap F:} Runtime error \#4 (substitute of unknown
        ingredient) had no test pinning the message.
  \item \textbf{Gap G:} Decision~\#31 (implicit Ingredient~$\to$~Quantity
        projection) had implementation coverage but no dedicated test
        class.
  \item \textbf{Gap I:} AmbiguousCall dispatcher had a fallback that
        dispatched via symbol-table isinstance checks -- silent-failure
        mode that would corrupt integration testing.
\end{itemize}

All three gaps were classified as ``fix tonight'' with concrete
remediation steps. I dispatched a second developer-agent round; that
round landed as \texttt{d5c539c} and the gaps were all closed. Test
count went 147 $\to$ 152.

\field{Rejected / modified.} Nothing rejected -- the three gaps were
all real.

\field{Errors caught.} The single most useful PM finding was Gap~I
(silent-AmbiguousCall fallback). It's the kind of bug that wouldn't
surface until İsmail's typechecker is fully wired in and tests are
running through the full pipeline -- at which point a missed
AmbiguousCall would corrupt the output. The PM caught it by reading
the interpreter source against the contract claim of ``loud raise on
AmbiguousCall.''

\field{Reflection.} Two PM rounds across both parts of this project
(entry~8 Part~1, entry~21 Part~2) have produced the same shape of
finding: silent-failure modes that would propagate downstream. The PM
doesn't catch obvious bugs -- the developer agent has tests for
those. The PM catches the kind of bug where two parts of the system
look fine in isolation but disagree under integration. \emph{Lesson:}
ask the PM to read contracts and implementations \emph{against each
other}, not separately.


\section{Session 13 -- Professor Agent: Pass-1 + Final Rubric Review}

Source: \texttt{session13\_professor\_2026-05-21\_to\_2026-05-23.jsonl}.
Same session UUID extended across two days: a Pass-1 narrow-scope review
on May~21 and a final full-scope review on May~23 after the type
checker and bug-fix rounds had landed. Two entries (22 and 23) cite
this one session.

\entryheader{22 of 24}{2026-05-21}{Pass-1 rubric review (autonomous agent)}{Claude Opus 4.7}

\field{Goal.} Get a tenured-PL-professor-style review of the
in-scope artifacts under the six-dimension rubric (Correctness, Design
Justification, Sebesta Framework, Originality, Code Quality, Exam
Performance) before submission, with predicted exam questions for each
partner.

\field{Prompt.} Professor-agent prompt (\textasciitilde 1.2k characters)
covering: scope (locked or stable artifacts only -- typechecker
explicitly out of scope at this point), required output structure
(rubric grades, top findings per dimension, predicted exam questions
with target answers and exposure ratings, triage list), tone
(skeptical, ``find what the grader will find'').

\field{Accepted.} Agent produced a structured $\sim$1500-word review
with: tentative letter grades per rubric dimension (B+ to A$-$ across
the six), 14 findings clustered by dimension, 5+5 predicted exam
questions with target answers and exposure ratings, and a triage list.

\field{Rejected / modified.} Nothing rejected substantively; every
finding I cross-checked against the actual code matched the agent's
claim.

\field{Errors caught.} F1 (the eager-vs-lazy step evaluation
inconsistency surfacing in sample~1's output) is the kind of finding I
would have walked into the exam with unprepared. The grader would have
asked ``why does sample~1 show serves~6 but only 4 loop iterations?''
F4 (error~\#19 referenced in plan~\S2.6 but missing from spec~\S10's
error list) and F5 (no Sebesta hook for why \texttt{scale} skips
Temperature/Duration) were both fixable in 5 minutes and not caught by
the PM review (which was looking for contract-vs-implementation drift,
not spec-vs-plan drift).

\field{Reflection.} The two-agent pattern I converged on by the end of
Part~2 -- PM agent for contract-vs-implementation drift, professor
agent for grader-vs-submission drift -- is the working pattern I'll
take into the next multi-week project. Each agent has a different
default failure mode: the PM defaults to ``this code passes its tests,
looks fine,'' the professor defaults to ``where would I subtract
points.''

\begin{tcolorbox}[colback=ackcolor!8,colframe=ackcolor!50,boxrule=0.5pt]
\small\textbf{Postscript (2026-05-23).} The F1 ``docstring fix + \S4.4
sentence'' resolution recorded above was superseded later that week.
After İsmail's bug-fix round introduced a \texttt{captured\_params}
re-execution mechanism for \texttt{scale}, I accepted the new
semantics (\texttt{scale} re-executes step bodies under the rebound
\texttt{servings}) rather than continuing to defend the eager-only
behavior with a docstring patch. The professor pass's F1 recommendation
became obsolete -- the better resolution was to take the free win the
new implementation offered. D1~\S4.4 and D3~\S2 were rewritten to
match the new semantics. The original docstring drift is now resolved
by code matching docs in the opposite direction from what the Pass-1
review proposed.
\end{tcolorbox}

\entryheader{23 of 24}{2026-05-23}{Final rubric review (autonomous agent, session 13 continuation)}{Claude Opus 4.7}

\field{Goal.} Final pre-submission rubric review of the complete Part~2
submission -- typechecker landed, bug-fix round integrated, samples
running end-to-end, documents drafted -- to confirm SHIP and surface
any last fixable findings.

\field{Prompt.} Same professor framing as Pass-1, scoped to the
finalized state: now includes the typechecker, the D1~[P2] sections,
the D3 PDF stub, D5, and the cumulative D4 in progress.

\field{Accepted.} Agent produced rubric grades, 5+5 predicted exam
questions per partner, three ranked fix-tonight items, and a clean
SHIP-WITH-FIXES verdict. Predicted exam questions for me cover the
\texttt{scale} re-execution rule (\textbf{B1}, low exposure -- defended
in D1~\S4.4 and D3~\S2); banker's rounding rationale (\textbf{B2}, low
exposure -- pinned by tests); Sebesta \S7.6 operand evaluation order
(\textbf{B3}, medium exposure -- rehearse the C-style contrast on
\texttt{i++ + i++}); the step-body partial-evaluator in
\texttt{interpreter.py} (\textbf{B4}, low exposure -- code is mine);
and Sebesta \S7.2 precedence of \texttt{quantity\_of} (\textbf{B5}, low
exposure -- decision \#34 rationale in spec). Predicted questions for
İsmail cover the type-equivalence split (\textbf{I1}, low),
\texttt{int}/\texttt{float} coercion asymmetry (\textbf{I2}, low),
the three-class \texttt{pour} verdict (\textbf{I3}, medium -- v2
refinement framing), \texttt{AmbiguousCall} rewrite-vs-annotate
(\textbf{I4}, medium -- must have read interpreter.py), and Pinch as
primitive (\textbf{I5}, low).

\field{Rejected / modified.} Nothing rejected. Three fix-tonight items
ranked by leverage: (F1)~edit one sentence in entry~22 to align the
journal narrative with the eventual fix path; (F3)~add error \#19 row
to spec~\S10's table; (F2)~wrap \texttt{env.define()} in
\texttt{\_eval\_LetStmt} with try/except to catch
\texttt{RedeclarationError}. All three are $\leq 10$~minutes each and
none touch tested code paths.

\field{Errors caught.} The F1 finding is the single most useful one
from this pass: my Pass-1 entry described one fix path (docstring +
\S4.4 sentence), and the eventual fix path was different (full
re-execution rule via \texttt{captured\_params}). Without the
postscript on entry~22 (now added), a grader who reads
D4~$\to$~D5~$\to$~code in order would catch the drift in 30 seconds.

\field{Reflection.} The same-session-extended-across-days pattern is
something I'll take forward. The professor agent had Pass-1 context
already in its working memory; the May~23 continuation didn't need
re-onboarding. The trade-off is that the cumulative output is harder
to grep against -- there are two reviews layered in one transcript --
but for personal use the continuity is worth it. Lesson for the
exam: \emph{the professor's job is to flag drift between artifacts;
the dev's job is to flag drift between code and tests; both jobs are
the human's job to integrate.}


\section{Session 14 -- Tester Agent: Final Pre-Submission Verification}

Source: \texttt{session14\_tester\_final\_verification\_2026-05-23.jsonl}.
Fresh session UUID; single long verification turn against the finalized
\texttt{origin/main} commit.

\entryheader{24 of 24}{2026-05-23}{Final functional verification (autonomous agent)}{Claude Opus 4.7}

\field{Goal.} Run the final pre-submission verification pass on the
Recipix Part~2 project against \texttt{origin/main}. Confirm 152~tests
green, all CLI commands behave as documented, sample-3 type-error
trace is pinned verbatim, the 19~\S10 errors are covered, and the
four canonical action-verb test cases produce the locked verdicts.

\field{Prompt.} Tester-agent prompt covering: scope (run, don't modify;
no commit, no push), required output structure (test count, CLI
results, sample-3 trace verbatim, 19-error coverage table, action-verb
verdicts, SHIP verdict), and explicit cleanup of any temporary
\texttt{.rcx} files created for the action-verb cases.

\field{Accepted.} Agent verified commit \texttt{db7a241}, ran all six
target CLI commands, built the 19-error coverage table, and ran the
four action-verb cases via temp \texttt{.rcx} files which it then
cleaned up. \textbf{Verdict: SHIP.} Concrete numbers:

\begin{itemize}[leftmargin=*,itemsep=2pt]
  \item \texttt{python -m unittest discover src/tests} $\to$ ``Ran
        152 tests in 0.039s -- OK'' (0 failures, 0 errors).
  \item \texttt{python main.py --run sample1.rcx} $\to$
        \texttt{pancakes -- serves 6} with the six pour/flip cycles
        (scale re-executes step bodies per \texttt{db7a241}).
  \item \texttt{python main.py --run sample2.rcx} $\to$ vegan smoothie
        with \texttt{milk} removed by substitute, then non-vegan with
        both \texttt{milk} and \texttt{oat\_milk} shown.
  \item \texttt{python main.py --typecheck sample3.rcx} $\to$ exit~1,
        \texttt{type error at line 13, col 0: dimension mismatch:
        cannot + Mass and Volume}.
  \item Imports smoke-test $\to$ \texttt{imports OK}.
\end{itemize}

The 19-error coverage table shows every error has a type-checker code
path with file:line citation. Action-verb verdicts: all four canonical
cases produce the expected accept/reject with the exact error messages.

\field{Rejected / modified.} Nothing rejected -- agent did exactly the
verification asked for, no scope creep into modifying code or
``helpful'' refactors. The two minor non-blocking observations the
tester surfaced -- (a)~no dedicated \texttt{test\_typechecker.py}
pinning per-error messages (regression fixtures provide coverage but
not regression-proof), and (b)~sample~3's column is reported as~0
rather than the column of \texttt{+} -- are both cosmetic.

\field{Errors caught.} None -- the submission was already clean before
this pass. The tester's contribution was \emph{confirming the SHIP
claim with concrete numbers}, which is the verification any
submission deserves before clicking upload. The verbatim sample-3
trace this session captured is the one pinned in D3~\S4.

\field{Reflection.} A final functional verification before submission
is cheap insurance. The tester agent's report adds nothing the test
suite doesn't already say, but the formatting -- exit codes, first 15
lines of output, the four action-verb temp programs run independently
of the existing test fixtures -- is exactly the artifact I want
sitting alongside the submission as proof-of-shippability. \emph{Lesson
for next time:} schedule the tester pass for the morning of the
deadline, not the night before. The tester catches nothing useful on
a hot tree where you're still editing; it catches everything useful on
a frozen tree where the only remaining action is \emph{click upload}.


\section*{Cumulative summary}

Across Part~1 (15 entries from 6 sessions, May~6--8) and Part~2 (this
continuation, 9 entries from 8 sessions, May~20--23), the journal
documents \textbf{14 substantive AI sessions producing 24 entries}.
The breakdown:

\begin{itemize}[leftmargin=*,itemsep=2pt]
  \item \textbf{6 Claude.ai web sessions} -- 2 brainstorming
        (Part~1 sessions 1 and 2), 1 Part~2 planning (session 7),
        3 E-experiments E1/E2/E3 (sessions 8, 9, 10).
  \item \textbf{2 Claude Code multi-turn PM threads} -- both in
        Part~1 (sessions 3 and 4 -- Sonnet 4.6 and Opus 4.7
        respectively, multi-turn user-driven orchestration).
  \item \textbf{6 Claude Code autonomous-agent runs} -- Part~1
        contributes 2 (tester agent session 5 + developer agent
        session 6 across two rounds); Part~2 contributes 4
        (developer session 11, PM session 12, professor session 13
        across two passes, tester session 14).
\end{itemize}

Mandatory experiments E1 (catch-up -- see entry~17 note), E2, and
E3 are all documented with shareable claude.ai/share URLs at
sessions~8, 9, 10. Four Claude Code transcripts are uploaded as
JSONL files alongside this PDF for sessions~11, 12, 13, 14.

The journal contains \textbf{seven documented cases} where the AI was
wrong, suboptimal, or made claims I had to verify against the spec,
the repo, or Sebesta and override:

\begin{itemize}[leftmargin=*,itemsep=2pt]
  \item Entry~2 (Part~1): structural equivalence suggested for
        \texttt{Recipe} type; caught by reading Sebesta~\S6.14
        against the suggestion.
  \item Entry~8 (Part~1): two parallel \texttt{tokens.py} files
        duck-typed compatibly but would silently break under rename;
        caught by the PM agent.
  \item Entry~14 (Part~1): developer agent generalized decision~\#17
        too broadly; caught by the tester-agent's
        \texttt{a < b == c} fixture.
  \item Entry~17 (Part~2 / E1): \texttt{substitute} slots typed as
        \texttt{Expression} instead of bare \texttt{IDENT}; caught by
        reading spec~\S13 against the AI's grammar.
  \item Entry~18 (Part~2 / E2): name equivalence recommended for
        \texttt{Ingredient} based on hypothetical record collisions
        that a closed v1 forecloses; caught by reading spec~\S1 and
        decision~\#10 against the recommendation.
  \item Entry~19 (Part~2 / E3): 12-row per-verb action-verb signature
        table instead of the 3-class collapse; caught by reading
        plan~\S2.3 against the AI's implementation.
  \item Entry~22/23 (Part~2 / Professor): F1 scale-vs-loop
        inconsistency in the Pass-1 review; the recommended fix
        (docstring patch) was superseded by a better fix
        (\texttt{captured\_params} re-execution rule) when İsmail's
        bug-fix round arrived. Caught by re-reading entry~22 in light
        of the May~22--23 changes.
\end{itemize}

That's seven documented AI-was-wrong moments across both parts of the
project -- above the ``three documented cases'' floor the handout
sets, and weighted toward the design / type-system corners where the
rubric's Design Justification and Sebesta Framework dimensions are
graded.

\vfill
\noindent\rule{\textwidth}{0.4pt}
\begin{center}
\small
\emph{End of Part~2 continuation. Combined with the Part~1 journal
(\texttt{210104004132\_D4\_P1.pdf}, May~8 submission), this completes
the D4 deliverable for CSE~341 Spring~2026.}
\end{center}

\end{document}
